\section{Background and Problem Formulation}

\subsection{RBF C-SVC dual}

For labels $y_i\in\{-1,+1\}$, the binary C-SVC dual is
\begin{equation}
  \min_{\bm{\alpha}}
  \quad
  \frac{1}{2}\bm{\alpha}^{\mathsf T}Q\bm{\alpha}
  -\bm{1}^{\mathsf T}\bm{\alpha},
  \label{eq:dual}
\end{equation}
subject to
\begin{equation}
  \bm{y}^{\mathsf T}\bm{\alpha}=0,
  \qquad
  0\leq \alpha_i\leq C.
  \label{eq:dual-constraints}
\end{equation}
The matrix entries are
\begin{equation}
  Q_{ij}=y_i y_j K(\bm{x}_i,\bm{x}_j),
  \qquad
  K(\bm{x}_i,\bm{x}_j)
  =\exp\!\left(-\gamma\lVert\bm{x}_i-\bm{x}_j\rVert_2^2\right).
  \label{eq:rbf}
\end{equation}
Materializing $Q\in\mathbb{R}^{N\times N}$ is prohibitive at scale.
Decomposition methods therefore optimize a working set while holding all
other variables fixed. SMO is the limiting two-variable method and uses an
analytic box-constrained update~\cite{platt1999smo}; second-order working-set
selection improves progress per step~\cite{fan2005working}.

\subsection{Execution pipeline}

The measured cuML path follows the five-stage pipeline in
\cref{fig:pipeline}.

\begin{figure}[t]
  \centering
  \begin{tikzpicture}[
      node distance=4.3mm,
      box/.style={draw=navy,rounded corners=2pt,fill=lightblue,
                  minimum width=0.78\textwidth,minimum height=7mm,
                  align=center,font=\small},
      hot/.style={draw=darkred,rounded corners=2pt,fill=red!6,
                  minimum width=0.78\textwidth,minimum height=7mm,
                  align=center,font=\small},
      arrow/.style={-{Stealth[length=2mm]},thick,color=darkgray}
    ]
    \node[box] (select) {Global working-set selection};
    \node[box,below=of select] (kww) {Materialize $K[W,W]$};
    \node[hot,below=of kww] (solve) {Local pair-SMO block solve\\
      \textbf{40.3\% of the steady solver phase}};
    \node[box,below=of solve] (global) {Generate or fetch $K[X,W]$ columns and update the global gradient};
    \node[box,below=of global] (stop) {Check the stopping condition and begin the next outer round};
    \draw[arrow] (select) -- (kww);
    \draw[arrow] (kww) -- (solve);
    \draw[arrow] (solve) -- (global);
    \draw[arrow] (global) -- (stop);
    \draw[arrow] (stop.east) -- ++(8mm,0) |- (select.east);
  \end{tikzpicture}
  \caption{Execution structure of the measured decomposition solver. The
  highlighted local solver becomes the dominant GPU kernel after coherent
  FP16 accelerates RBF materialization and cache traffic.}
  \label{fig:pipeline}
\end{figure}

The working-set size is typically 1,024 or 4,096. Two kernel-data units must
be distinguished:

\begin{itemize}
  \item a \emph{global column} $K[X,j]$ contains $N$ RBF values and requires
        scanning the full $N\times d$ feature matrix on a cache miss; and
  \item a \emph{local row} $K[W,j]$ contains $|W|$ values already present in
        the local $K[W,W]$ matrix and is consumed by the block solver.
\end{itemize}

They have different units and costs and must not be added as raw counts.

\subsection{Local working-set QP and serial dependence}

Define the signed working-set displacement
\begin{equation}
  s_i=y_i(\alpha_i-\alpha_i^0).
\end{equation}
The captured local problem is
\begin{equation}
  \min_{\bm{s}}
  \quad
  \bm{f}_0^{\mathsf T}\bm{s}
  +\frac{1}{2}\bm{s}^{\mathsf T}K\bm{s},
  \label{eq:local-qp}
\end{equation}
subject to
\begin{equation}
  \bm{1}^{\mathsf T}\bm{s}=0,
  \qquad
  \bm{\ell}\leq\bm{s}\leq\bm{u}.
\end{equation}
Each pair update selects an upper variable $u_t$ and a lower variable $l_t$,
computes a feasible step $q_t$, and updates the entire local gradient:
\begin{equation}
  \bm{f}^{t+1}
  =\bm{f}^{t}
  +q_t\left(K_{u_t,:}-K_{l_t,:}\right).
  \label{eq:pair-update}
\end{equation}
The next pair is selected from $\bm{f}^{t+1}$:
\begin{equation}
  (u_t,l_t)\rightarrow q_t\rightarrow \bm{f}^{t+1}
  \rightarrow (u_{t+1},l_{t+1}).
  \label{eq:serial-chain}
\end{equation}
For a valid positive-semidefinite kernel, this is a convex QP rather than an
NP-hard problem. The systems challenge is to reduce dependency depth without
increasing dense-algebra work, local-row consumption, global-column misses, or
outer rounds.

\subsection{Cost model}

We model the solver cost by three non-interchangeable terms:
\begin{equation}
  T \approx
  M_g\,C_g(N,d)
  + R_\ell\,C_\ell(|W|)
  + D\,C_d
  + T_{\mathrm{misc}},
  \label{eq:cost-model}
\end{equation}
where $M_g$ is the number of global $K[X,j]$ cache misses, $C_g$ is the cost
of producing one such column, $R_\ell$ is the number of local $K[W,j]$ row
uses, and $D$ is the number of dependent decision/update stages. A useful
optimization must reduce at least one term without causing a larger increase
in another. This model is the basis of the solver gates in
\cref{sec:boundaries}.
